\documentclass[11pt]{letter}
\usepackage[margin=1in]{geometry}
\usepackage{setspace}
\setstretch{1.05}

\signature{Yuguo Yu, on behalf of all authors\\
Institute of Science and Technology for Brain-Inspired Intelligence\\
Fudan University, Shanghai, China\\
\texttt{yuyuguo@fudan.edu.cn}}

\address{Mingyi Huang, Wei Lin, Anna Wang Roe, Yuguo Yu\\
Fudan University \& Zhejiang University\\
\\
[DATE OF SUBMISSION]}

\begin{document}

\begin{letter}{The Editors\\Nature Communications}

\opening{Dear Editor,}

We are pleased to submit our manuscript \textit{``Gibbs Manifold Geometry Unifies Neural Variability, Sparseness, and Metabolic Efficiency''} for consideration as an Article in \textit{Nature Communications}.

\textbf{Specific advance.}
Cortical neurons fire variably, sparsely, and under a tight metabolic budget---three features long studied in isolation by sparse-coding, balanced-network, and efficient-coding theories.  We construct, to our knowledge for the first time, a \textit{Gibbs manifold}---the one-parameter family of optimal response distributions arising from entropy maximization under a metabolic budget---and show analytically that the coefficient of variation (CV) is the macroscopic coordinate along it that simultaneously controls sparseness, information--cost efficiency, and (under Poisson-like statistics) trial-to-trial reliability.  The framework distinguishes three logical statuses---algebraic identity, normative prediction, and model-dependent corollary---making each claim precise and falsifiable.

\textbf{Headline results.}
(i) An FR-controlled per-neuron decoder analysis on macaque V1 ($N=564$ neurons, 5 sessions) shows lifetime CV predicts decodable stimulus information beyond mean firing rate (partial Spearman $\rho_{|\,\mathrm{FR}}=+0.68$, $p=10^{-78}$; Fisher-combined $\rho=+0.72$, $p=10^{-100}$), converting the theoretical CV--reliability link into an empirical, FR-controlled relationship.
(ii) Replacing hand-tuned sparsity priors with population-CV maximization recovers V1-like Gabor receptive fields and yields a more favourable efficiency--fidelity Pareto position than the Cauchy-prior SparseNet baseline---suggesting CV as a hardware-agnostic objective for neuromorphic representations.
(iii) V1 data, single-neuron simulations, and recurrent E--I networks trace concave trajectories consistent with the Gibbs manifold on a common entropy--cost plane after system-specific rate calibration.

\textbf{Broad readership.}
The work bridges computational neuroscience (efficient coding, balanced networks, E--I plasticity), information theory (constrained entropy maximization, Pareto frontiers), statistical physics (the Gibbs/Helmholtz analogue), and machine learning (sparse coding, energy-aware representation learning).

\textbf{Honest scope.}
The strictly-convex-cost theorem covers $\mathrm{CV}\in[1/\sqrt{3},1]$; empirical $\mathrm{CV}>1$ in cortex is accommodated by a phenomenological sublinear-cost extension whose monotonicity we verify numerically rather than analytically.  The deterministic-encoding approximation $I\approx H$ does not hold strictly on our V1 dataset ($\bar\rho\approx 0.85$); the manifold then serves as a within-dataset upper bound on attainable $I/C$, discussed and cross-checked across datasets in the Supplement.

\textbf{Code, data, reporting, LLM use.}
Custom code for all analyses, simulations, and figures is available to editors and reviewers via a private GitHub repository (\textbf{[private review URL]}).  Source-data Excel files for every figure panel, the Nature Portfolio Reporting Summary, and the Supplementary Information are uploaded with this submission.  V1 data are publicly available via CRCNS (\texttt{v1v2-1}).  Language polishing of an earlier draft was assisted by an LLM (Claude, Anthropic), as declared in the Acknowledgements; all scientific content, derivations, code, and conclusions are the authors' own.

\textbf{Suggested reviewers}, none with recent collaboration or perceived conflict:
\begin{enumerate}
\itemsep -3pt
\item \textbf{[NAME]} ([Affiliation]) --- efficient coding / information theory in V1; \texttt{[email]}
\item \textbf{[NAME]} ([Affiliation]) --- balanced E--I networks and AI cortical dynamics; \texttt{[email]}
\item \textbf{[NAME]} ([Affiliation]) --- cortical metabolism and energy-efficient coding; \texttt{[email]}
\item \textbf{[NAME]} ([Affiliation]) --- maximum-entropy models of neural populations; \texttt{[email]}
\item \textbf{[NAME]} ([Affiliation]) --- neuromorphic / energy-efficient learning; \texttt{[email]}
\end{enumerate}
\textbf{Excluded reviewers:} \textbf{[NAME(S), with one-line reason --- or ``no exclusion requests'']}.

\textbf{Prior discussion / submission history.}
We have had no prior pre-submission discussion with a \textit{Nature Communications} editor about this work.  \textbf{[If applicable: ``An earlier version was evaluated at [JOURNAL] and not sent to review; we have substantially revised the framework, added the multi-session FR-controlled decoder analysis, tightened the convex- vs.\ sublinear-cost boundary, and narrowed the central claims.''  Delete if not applicable.]}

The manuscript has not been published, and is not under consideration for publication, elsewhere.  All authors have read and approved the submitted version.

\closing{Sincerely,}

\end{letter}

\end{document}
